%=======================================================================
%  Paper V — Section 6 (The dividing plane is not self-determined) FINAL
%
%  Requires the macros listed at the head of Section 2.
%
%  PRINCIPAL CHANGE IN THIS REVISION ---------------------------------
%  The argument no longer rests on the extrapolated turning point. The
%  turning point quoted in ver10 (0.253-0.255 km from cubic and quartic
%  fits) falls within a few metres of the plane at which the 8000 kg/m3
%  scan ceiling of Section 2.9 begins to truncate the recovered
%  contrast (x ~ 0.257 km), and local extrapolation of the gradient
%  puts it at 0.262-0.278 km instead. It is therefore demoted to a
%  remark with its caveats stated. The argument is instead closed on
%  observed data alone: the improvement is monotone across all thirteen
%  planes, and the outermost plane exceeds the grain-density ceiling by
%  11.2 sigma, so the maximum over the admissible set is attained at a
%  boundary fixed by composition rather than by the objective.
%
%  Superseded values:
%      sigma quoted to 6 dp      -> 5 dp (grid floor is 3.5e-6)
%      I quoted to 3 dp          -> 2 dp
%      "a turning point must exist" -> not proved; limit at the tip is
%                                   finite and the minimum may be at the
%                                   endpoint
%      "genuine convergence"     -> consistency with a one-sided bound
%      "7.5 kg/m3 per metre"     -> 10.0 at the reference plane; and the
%                                   sensitivity of rho_rest, 0.28, is
%                                   the one that matters
%      eq:massconstraint         -> eq:massbalance
%
%  All values confirmed against paperV_final.json. The rate-of-fall list
%  is reconciled below: twelve gradients, not thirteen.
%=======================================================================

\section{The dividing plane is not self-determined}
\label{sec:plane}

The dividing plane is the one ingredient of the method that is neither
measured nor derived. The mass comes from radio science, the shape from a
spacecraft, the spin from lightcurves and imaging; the plane is placed by
the investigator. Section~\ref{sec:itokawa:sweep} showed that the recovered
region density depends strongly on where it is placed. This section asks
whether the objective function can place it, and finds that it cannot---not
in the weak sense that the optimum is shallow, but in the strong sense that
within the range the composition permits there is no interior optimum at
all, and the objective's preferred plane is one the meteorite collection
forbids.

\subsection{The improvement rises monotonically to the end of the sweep}
\label{sec:plane:monotonic}

Across the whole Itokawa sweep the dispersion at the minimum falls
monotonically, from $0.07549$ at $\xsplit = \SI{0.130}{\kilo\metre}$ to
$0.06373$ at $\SI{0.240}{\kilo\metre}$, against a uniform-body value of
$0.08663$: an improvement rising from $12.86\%$ to $26.44\%$ without a
single reversal in thirteen planes (Table~\ref{tab:itokawasweep},
Figure~\ref{fig:plane}).

We quote the dispersion to five decimal places and the improvement to two,
and not further, although the computation is deterministic and returns more
digits. The reason is the scan grid of
Section~\ref{sec:methods:minimisation}. Near its minimum the objective is
quadratic, $\potvarn(\densR) \simeq \potvarn^{\ast} +
\tfrac{1}{2}k(\densR-\densR^{\ast})^{2}$, and the curvature sensitivity
interval---a rise of $0.60\%$ over
$\pm\SI{130}{\kilo\gram\per\cubic\metre}$---fixes
$k = 7.1\times10^{-7}\,\potvarn^{\ast}$ per $(\mathrm{kg\,m^{-3}})^{2}$.
A half-step of $\SI{12.5}{\kilo\gram\per\cubic\metre}$ therefore inflates
the reported minimum by at most $5.5\times10^{-5}$ in relative terms, or
$3.5\times10^{-6}$ absolute. That is below the step-to-step change along the
sweep, roughly $1.1\times10^{-4}$ per ten metres, so the monotonicity is
safe; but it is above the sixth decimal place, so the sixth decimal place is
grid noise and we do not print it.

The rate of fall is not uniform. It steepens from about
$-1.11\times10^{-3}$ per ten metres at the start of the sweep to a maximum
near $\SI{0.185}{\kilo\metre}$, then slackens to $-1.06\times10^{-3}$
between $0.200$ and $\SI{0.220}{\kilo\metre}$ and to
$-6.9\times10^{-4}$ between $0.220$ and $\SI{0.240}{\kilo\metre}$. The
slackening is the first sign that the fall is not indefinite, and
Section~\ref{sec:plane:impossible} takes it up.
% Resolved: there is no fourteenth plane. The thirteen-entry list came from
% a run record in which x = 0.1608 km appears twice, the reference run and
% the sweep run at the same detected plane, differing only in the sixth
% digit of the excess mass. De-duplicated, the thirteen planes give twelve
% gradients, in units of 1e-6 per ten metres:
%   -1107 -1130 -1163 -1178 -1196 -1222 -1237 -1230 -1223 -1194 -1061 -692
% The same duplicate explains the two earlier discrepancies noted in the
% working file; the sweep driver emitted the plane set Table 3.1 prints.

\begin{figure}[htbp]
\centering
\includegraphics[width=\textwidth]{fig08_improvement_vs_plane.pdf}
\caption{Dispersion improvement against dividing-plane position on
Itokawa. The improvement rises monotonically across all thirteen planes
with no interior turning point. The vertical bands mark the two external
bounds of Section~\ref{sec:itokawa:ceiling}: the empirical block-density
bound at $\SI{194+-11}{\metre}$ and the LL-chondrite grain-density ceiling
at $\SI{207+-5}{\metre}$. The plane returned by the neck detector, at
$\SI{160.8}{\metre}$, is marked; the objective prefers the block-density
boundary to it by $4.6$ percentage points of improvement. The shaded
horizontal band is the $0.9$--$5.9\%$ mesh-resolution envelope of
Paper~IV, shown on the same scale to indicate that the objective's
preference among admissible planes is of the same order as its own
resolution sensitivity.}
\label{fig:plane}
\end{figure}

\subsection{Why this settles the question without extrapolation}
\label{sec:plane:impossible}

The monotonicity by itself does not show that the objective cannot choose a
plane; a function rising to the edge of the range examined might turn just
beyond it. What settles the question is that the planes the objective
prefers are excluded on grounds that have nothing to do with the objective.

At $\xsplit = \SI{0.240}{\kilo\metre}$ the recovered head density is
$\SI{5000}{\kilo\gram\per\cubic\metre}$. The LL-chondrite grain density is
$3540 \pm \SI{130}{\kilo\gram\per\cubic\metre}$
(Section~\ref{sec:itokawa:ceiling}), so the objective's best available plane
sits $11.2$ standard deviations above the density of solid LL material with
no pore space at all. It is not a marginal exclusion and it does not depend
on which of the two bounds is adopted: the same plane exceeds the empirical
block-density bound of $3220 \pm 220$ by $8.1$ standard deviations.

The argument then closes on observed data alone. The improvement is strictly
increasing over every plane in the sweep; the admissible set is the interval
below $\xsplit = \SI{194+-11}{\metre}$; therefore the improvement attains its
maximum over the admissible set at the upper boundary of that set, and the
boundary is fixed by the meteorite collection rather than by the
minimization. No statement about the behaviour of $\potvarn$ beyond the
sweep is required.

It is worth putting a number on what the objective's preference costs.
Interpolating Table~\ref{tab:itokawasweep} to the block-density boundary
gives $\improve \simeq 21.5\%$ there, against $16.89\%$ at the neck the
detector returns: the objective prefers the boundary by $4.6$ percentage
points, or $27\%$ in relative terms. That difference is real, but it lies
inside the $0.9$--$5.9\%$ mesh-resolution envelope of Paper~IV, so it is not
a discrimination the method can defend on its own terms. Over the full
admissible range the improvement varies by $8.6$ percentage points, which
does exceed the envelope; the objective discriminates among admissible
planes, but only in the direction of the boundary.

\paragraph{The turning point, and why we do not use it.}
A cubic fit to $\potvarn(\xsplit)$ over the sweep places a minimum at
$\SI{0.253}{\kilo\metre}$ and a quartic fit at $\SI{0.255}{\kilo\metre}$.
Earlier drafts of this section reported those figures as the location of the
turning point. We now report them only as a remark, for three reasons, and
the argument above is constructed so as not to depend on them.

First, they are extrapolations of $13$~m to $15$~m beyond the last plane
computed, and the agreement between the two fit orders is not evidence of
accuracy: two adjacent polynomial orders extrapolate similarly by
construction. Extrapolating the measured gradient instead places the zero at
$\SI{0.262}{\kilo\metre}$ on a quadratic through the last three planes and
$\SI{0.278}{\kilo\metre}$ on a straight line through the last two, which is
$8$ to $25$~m further out.

Second, the extrapolated position coincides with the plane at which the scan
ceiling begins to bind. Holding $\Dmass$ at $2.80\times10^{9}$~kg and
$\densrest$ near $\SI{1855}{\kilo\gram\per\cubic\metre}$, the ceiling of
$\SI{8000}{\kilo\gram\per\cubic\metre}$ used at the outer planes
(Section~\ref{sec:itokawa:sweep}) permits a region no smaller than
%
\begin{equation}
f_{\min} = \frac{\Dmass}{(\densR^{\max}-\densrest)\,\Vtot}
         = \frac{2.80\times10^{9}}
                {6145 \times 1.77856\times10^{7}}
         = 2.56\% ,
\label{eq:capbinds}
\end{equation}
%
which the sweep reaches at $\xsplit \approx \SI{0.257}{\kilo\metre}$. A
turning point predicted at $0.253$ to $\SI{0.255}{\kilo\metre}$ therefore
lies within a few metres of the plane beyond which the recovered contrast
would be truncated by the scan rather than determined by the data, and the
two cannot be separated in the present configuration. Re-running the outer
planes with the ceiling raised to, say,
$\SI{20000}{\kilo\gram\per\cubic\metre}$ would separate them; we have not
done so, because the result would lie far outside the compositional bounds
either way.

Third, we do not assert that a turning point must exist at all. As the plane
advances toward the tip, $\VR \to 0$ while the excess mass need not: the
region approaches a point mass of free magnitude $\Dmass$ located at the
tip, and the dispersion approaches the finite value obtained by optimizing
over that magnitude. The objective is therefore continuous on a closed
interval with a finite limit at its right endpoint, and its minimum may be
attained at that endpoint rather than at an interior stationary point.
Whether it is attained interior or at the limit is a question the present
data do not answer and the argument of this section does not need answered.

\subsection{Why this happens}
\label{sec:plane:why}

The explanation is Section~\ref{sec:constrains}. What the minimization
determines in the strong-signal regime is the background density, and with
it, through Equation~\eqref{eq:excessidentity}, the excess mass. Write the
region density as the background density plus the excess mass spread over
the region,
%
\begin{equation}
\densR = \densrest + \frac{\Dmass}{\VR} ,
\label{eq:plane:decompose}
\end{equation}
%
which is Equation~\eqref{eq:excessidentity} rearranged. As the plane
advances, $\VR$ falls, the mass constraint~\eqref{eq:massbalance} carries
$\densrest$ toward $\Mtot/\Vtot$, and the objective compensates by raising
$\densR$. Because the perturbation to the surface potential enters
Equation~\eqref{eq:superposition} through the product
$(\densR-\densrest)\UR$, and $\UR$ scales with $\VR$ for a small region, the
objective can hold that product nearly fixed while $\VR$ shrinks. It has no
reason to stop, and within the range examined it does not.

This is also the reconciliation between this section and
Section~\ref{sec:constrains}, which at first reading point in opposite
directions: one says the plane is unconstrained, the other that the answer
is stable under moving it. Both are true, and they are compatible precisely
because the stable quantity is not the one being optimized. Moving the plane
from the detected neck to the block-density boundary---a displacement of
$\SI{32.7}{\metre}$, a fifth of the plane's own position---changes the
recovered quantities as follows:
%
\begin{multline}
\denshead:\ 2800 \to 3220\ (+15.0\%),
\qquad
\Dmass:\ 2.753 \to 2.881\ (+4.65\%),
\qquad \\
\densrest:\ 1858.1 \to 1850.9\ (-0.387\%).
\label{eq:plane:move}
\end{multline}
%
The ratio of the last two is $4.65/0.387 = 12.0$, which is the amplification
factor $\ampfac$ at the reference plane, recovered here from a finite
displacement rather than from the differential
Equation~\eqref{eq:excessamplify}; the ratio of the first and last is
$38.8$, or $\ampfac$ times the residual factor of $3.2$ that measures
compensation between contrast and region volume. A plane displacement large
enough to move the region density by fifteen per cent moves the quantity the
method constrains by four parts in a thousand.

\subsection{What must supply the plane instead}
\label{sec:plane:external}

Two external criteria are available, and on Itokawa they do not contradict
one another. We put it that way deliberately, because an earlier draft of
this section described their relation as a convergence, and that is more
than the evidence supports.

The first is geomorphological. The neck detector of
Section~\ref{sec:methods:neck} returns
$\xsplit = \SI{160.8}{\metre}$ at a prominence of $0.170$, using only the
cross-sectional profile of the shape model and no density information
whatever. The second is compositional: the block-density bound of
Section~\ref{sec:itokawa:ceiling} excludes planes beyond
$\SI{194+-11}{\metre}$, and the grain-density ceiling those beyond
$\SI{207+-5}{\metre}$, using meteorite measurements and no shape information
beyond the gradient used to convert a density into a plane position.

The two are genuinely independent in their inputs, and the detected neck
falls $\SI{33}{\metre}$ inside the tighter of the two bounds, about three
times that bound's own uncertainty. But the compositional criterion is
one-sided: it excludes planes that are too far out and says nothing about
planes that are too far in. Every plane in the sweep below
$\SI{194}{\metre}$ satisfies it, which is $63.5$ of the $110$~m swept, so
under a uniform null a plane chosen at random would have satisfied it $58\%$
of the time. The correct statement is that an independently chosen plane
turns out not to be excluded, which is reassurance rather than corroboration,
and not that two methods have selected the same plane. We adopt the detected
neck because it has geometric warrant, and we use the compositional bound to
truncate the sweep, not to locate it.

One apparent inconsistency should be disposed of. The neck slab of
Section~\ref{sec:itokawa:ambiguity} has its upper edge at
$\SI{202.9}{\metre}$, beyond the block-density bound of
$\SI{194}{\metre}$. There is no conflict: the bound is a bound on the
recovered \emph{density}, converted into a plane position by
$\mathrm{d}\denshead/\mathrm{d}\xsplit$ for the \emph{head} model, and the
conversion does not transfer to a different region shape. The neck model
returns $\densneck = \SI{2525}{\kilo\gram\per\cubic\metre}$, comfortably
below both bounds, so it violates neither.

The geomorphological criterion carries its own resolution dependence, and it
propagates asymmetrically. Section~\ref{sec:itokawa:mesh} found that
decimating to $20{,}000$ facets displaces the detected plane by
$\SI{9.4}{\metre}$. At the local gradients of
Table~\ref{tab:itokawasweep}---$\SI{10.0}{\kilo\gram\per\cubic\metre\per
\metre}$ for $\denshead$ and
$-\SI{0.28}{\kilo\gram\per\cubic\metre\per\metre}$ for $\densrest$---that
displacement moves the region density by $\SI{94}{\kilo\gram\per
\cubic\metre}$, or $3.4\%$, which is comparable to the curvature sensitivity
interval of $4.6\%$; and the background density by
$\SI{2.6}{\kilo\gram\per\cubic\metre}$, or $0.14\%$, against a combined
budget of $1.45\%$, where it adds nothing in quadrature. The plane detector
is a significant error source for the quantity the method does not constrain
and a negligible one for the quantity it does.

We note that this is a weaker footing than the rest of the method. The
resolution stability of the inferred density
(Section~\ref{sec:itokawa:mesh}) and the invariance of the background
density in the strong-signal regime (Section~\ref{sec:constrains}) are
properties measured within the method. The placement of the plane is not,
and rests on judgement informed by two independent lines of evidence, only
one of which is two-sided. Any result reported from this method should state
the plane used, the criterion by which it was chosen, and the sensitivity of
the answer to moving it---and should state that sensitivity for the quantity
being reported. On Itokawa at the reference plane the three sensitivities
are $\SI{10.0}{\kilo\gram\per\cubic\metre\per\metre}$ for $\denshead$
(rising to $\SI{55}{\kilo\gram\per\cubic\metre\per\metre}$ at the outer
planes and falling to $7.5$ over the first four),
$\SI{0.28}{\kilo\gram\per\cubic\metre\per\metre}$ for $\densrest$, and
$0.18\%$ per metre for $\Dmass$. Quoting the first alone, as earlier drafts
did, states the sensitivity of the quantity
Section~\ref{sec:constrains} concludes the method does not determine.